\sectionheading{5.23 অ্যালগরিদম (Algorithm)}

কোনো Problem সমাধানের জন্য সরাসরি Programming Language-এ
Code লেখা শুরু না করে প্রথমে সমাধানের কাজগুলো
\textbf{কোন ক্রমে সম্পন্ন হবে}, তা ধাপে ধাপে সাজিয়ে নেওয়া ভালো।

Problem সমাধানের এই ধারাবাহিক ও যৌক্তিক ধাপসমূহই হলো
\textbf{Algorithm}।

%================================================
% Definition
%================================================
\booksubheading{Algorithm কাকে বলে?}

\begin{definitionbox}

\textbf{\color{darkTeal}সংজ্ঞা}

কোনো নির্দিষ্ট Problem সমাধানের জন্য প্রয়োজনীয় কাজগুলোকে
\textbf{যুক্তিসংগত ও পর্যায়ক্রমিক ধাপে সাজিয়ে লিখিতভাবে প্রকাশ
করাকে Algorithm} বলে।

\end{definitionbox}

\textbf{সহজভাবে মনে রাখবে—}

\begin{center}

\textbf{
Problem
\(\longrightarrow\)
ধাপে ধাপে Solution
\(\longrightarrow\)
Algorithm
}

\end{center}

অর্থাৎ Algorithm বলে দেয়—
\textbf{কোন কাজটি আগে হবে, তার পরে কী হবে এবং কীভাবে কাঙ্ক্ষিত
Result পাওয়া যাবে।}

%================================================
% Origin
%================================================
\booksubheading{Algorithm শব্দের উৎপত্তি — সংক্ষিপ্ত তথ্য}

Algorithm শব্দটির উৎপত্তি গণিতবিদ
\textbf{মুসা আল-খারিজমি (Musa Al-Khwarizmi)}
এর নামের সঙ্গে সম্পর্কিত।

\begin{rememberbox}

\textbf{\color{deepOrange}MCQ তথ্য}

\textbf{Algorithm} শব্দটির সঙ্গে
\textbf{Musa Al-Khwarizmi}-এর নাম সম্পর্কিত।

এটি মূল ধারণার তুলনায় কম গুরুত্বপূর্ণ;
MCQ-এর জন্য সংক্ষেপে মনে রাখলেই যথেষ্ট।

\end{rememberbox}

%================================================
% Importance
%================================================
\booksubheading{Algorithm কেন প্রয়োজন?}

Program লেখার আগে Algorithm তৈরি করলে—

\begin{compacttwocolpoints}

\item Problem-এর Solution সহজে বোঝা যায়।

\item কাজের সঠিক ক্রম নির্ধারণ করা যায়।

\item Programming Language-এ Code লেখা সহজ হয়।

\item Program-এর Logic আগেই যাচাই করা যায়।

\item ভুল শনাক্ত ও সংশোধনে সুবিধা হয়।

\item অন্য কেউ Solution সহজে বুঝতে পারে।

\item পরবর্তীতে Program পরিবর্তন বা উন্নয়ন করা সহজ হয়।

\end{compacttwocolpoints}

%================================================
% Rules
%================================================
\booksubheading{Algorithm লেখার সাধারণ নিয়ম}

Algorithm লেখার সময়—

\begin{enumerate}

\item প্রথম ধাপে \textbf{শুরু} লিখতে হয়।

\item প্রয়োজনীয় \textbf{Input} গ্রহণের ধাপ লিখতে হয়।

\item Calculation, Formula বা প্রয়োজনীয় কাজ সঠিক ক্রমে লিখতে হয়।

\item প্রয়োজনীয় \textbf{Output} প্রদর্শনের ধাপ লিখতে হয়।

\item সর্বশেষ ধাপে \textbf{শেষ} লিখতে হয়।

\item প্রতিটি ধাপ সংক্ষিপ্ত, স্পষ্ট ও দ্ব্যর্থহীন হওয়া উচিত।

\end{enumerate}

\textbf{সাধারণ কাঠামো—}

\begin{center}
\begin{adjustbox}{max width=0.98\linewidth}
\fcolorbox{darkBlue}{softBlue}{
    \parbox[c][8mm][c]{0.13\linewidth}{
        \centering\textbf{শুরু}
    }
}
\(\longrightarrow\)
\fcolorbox{mainTeal}{softGreen}{
    \parbox[c][8mm][c]{0.14\linewidth}{
        \centering\textbf{Input}
    }
}
\(\longrightarrow\)
\fcolorbox{darkBlue}{softBlue}{
    \parbox[c][8mm][c]{0.15\linewidth}{
        \centering\textbf{Process}
    }
}
\(\longrightarrow\)
\fcolorbox{mainTeal}{softGreen}{
    \parbox[c][8mm][c]{0.14\linewidth}{
        \centering\textbf{Output}
    }
}
\(\longrightarrow\)
\fcolorbox{darkBlue}{softBlue}{
    \parbox[c][8mm][c]{0.13\linewidth}{
        \centering\textbf{শেষ}
    }
}
\end{adjustbox}
\end{center}

%================================================
% Examples
%================================================
\booksubheading{Algorithm-এর উদাহরণ}

%================================================
% Example 1
%================================================
\begin{examplebox}{উদাহরণ–১: দুটি সংখ্যার যোগফল নির্ণয়}

\textbf{Problem:}
দুটি সংখ্যার যোগফল নির্ণয় করতে হবে।

\medskip

\textbf{\color{darkTeal}Algorithm}

\textbf{ধাপ–১:} শুরু।

\textbf{ধাপ–২:} \textenglish{\ttfamily a} ও
\textenglish{\ttfamily b} দুটি সংখ্যা গ্রহণ করি।

\textbf{ধাপ–৩:}
\textenglish{\ttfamily sum = a + b} নির্ণয় করি।

\textbf{ধাপ–৪:}
\textenglish{\ttfamily sum} প্রদর্শন করি।

\textbf{ধাপ–৫:} শেষ।

\medskip

\textbf{সংক্ষেপে—}

\begin{center}

\textbf{
Input a, b
\(\longrightarrow\)
sum = a + b
\(\longrightarrow\)
Print sum
}

\end{center}

\textbf{\color{darkTeal}Algorithm থেকে Python Code}

\begin{codebox}
\begin{english}
\ttfamily
a = int(input("Enter first number: "))\\
b = int(input("Enter second number: "))\\[0.8mm]
sum = a + b\\[0.8mm]
print("Sum =", sum)
\end{english}
\end{codebox}

এখানে—

\begin{itemize}

\item Algorithm-এর \textbf{Input}
\(\rightarrow\) Python-এর \textenglish{\ttfamily input()}

\item Algorithm-এর \textbf{Process}
\(\rightarrow\) \textenglish{\ttfamily sum = a + b}

\item Algorithm-এর \textbf{Output}
\(\rightarrow\) \textenglish{\ttfamily print()}

\end{itemize}

\end{examplebox}

\begin{rememberbox}

\textbf{\color{deepOrange}মনে রাখবে}

এখানে Python Code-এর উদ্দেশ্য শুধু
\textbf{Algorithm কীভাবে Program Code-এ রূপ নেয়} তা বোঝানো।

Python-এর Input/Output পরবর্তী অংশে বিস্তারিত আলোচনা করা হবে।

\end{rememberbox}

%================================================
% Example 2
%================================================
\begin{examplebox}{উদাহরণ–২: ত্রিভুজের ক্ষেত্রফল নির্ণয়}

\textbf{Problem:}
ত্রিভুজের ভূমি ও উচ্চতা দেওয়া থাকলে ক্ষেত্রফল নির্ণয় করতে হবে।

\medskip

\textbf{Formula:}

\begin{center}

\textbf{
Area = \(\frac{1}{2}\) \(\times\) ভূমি \(\times\) উচ্চতা
}

\end{center}

\textbf{\color{darkTeal}Algorithm}

\textbf{ধাপ–১:} শুরু।

\textbf{ধাপ–২:} ভূমি \textenglish{\ttfamily b} এবং
উচ্চতা \textenglish{\ttfamily h} গ্রহণ করি।

\textbf{ধাপ–৩:}
\textenglish{\ttfamily Area = 1/2 * b * h} হিসাব করি।

\textbf{ধাপ–৪:}
\textenglish{\ttfamily Area} প্রদর্শন করি।

\textbf{ধাপ–৫:} শেষ।

\end{examplebox}

%================================================
% Example 3
%================================================
\begin{examplebox}{উদাহরণ–৩: পাঁচটি সংখ্যার যোগফল নির্ণয়}

\textbf{Problem:}
পাঁচটি সংখ্যার যোগফল নির্ণয় করতে হবে।

\medskip

\textbf{\color{darkTeal}Algorithm}

\textbf{ধাপ–১:} শুরু।

\textbf{ধাপ–২:}
\textenglish{\ttfamily A, B, C, D, E}
পাঁচটি সংখ্যা গ্রহণ করি।

\textbf{ধাপ–৩:}

\begin{center}

\textenglish{\ttfamily S = A + B + C + D + E}
\end{center}

নির্ণয় করি।

\textbf{ধাপ–৪:}
\textenglish{\ttfamily S} প্রদর্শন করি।

\textbf{ধাপ–৫:} শেষ।

\medskip

\textbf{Shortcut—}

\begin{center}

\textbf{
A, B, C, D, E
\(\longrightarrow\)
যোগ
\(\longrightarrow\)
S
}

\end{center}

\end{examplebox}

%================================================
% Example 4
%================================================
\begin{examplebox}{উদাহরণ–৪: তিনটি সংখ্যার গড় নির্ণয়}

\textbf{Problem:}
তিনটি সংখ্যার Average নির্ণয় করতে হবে।

\medskip

\textbf{\color{darkTeal}Algorithm}

\textbf{ধাপ–১:} শুরু।

\textbf{ধাপ–২:}
\textenglish{\ttfamily a, b, c}
তিনটি সংখ্যা গ্রহণ করি।

\textbf{ধাপ–৩:}

\begin{center}

\textenglish{\ttfamily sum = a + b + c}
\end{center}

নির্ণয় করি।

\textbf{ধাপ–৪:}

\begin{center}

\textenglish{\ttfamily avg = sum / 3}
\end{center}

নির্ণয় করি।

\textbf{ধাপ–৫:}
\textenglish{\ttfamily avg} প্রদর্শন করি।

\textbf{ধাপ–৬:} শেষ।

\medskip

এখানে একটি Problem সমাধানের জন্য
\textbf{একাধিক Processing Step} ব্যবহার করা হয়েছে।

\end{examplebox}

%================================================
% Types of Logic
%================================================
\booksubheading{Algorithm-এর সাধারণ Logic সহজভাবে}

Algorithm-এর Solution বিভিন্ন ধরনের Logic ব্যবহার করে তৈরি হতে পারে।
মূল ধারণাগুলো হলো—

\begin{tabularx}{\linewidth}{|>{\bfseries}p{0.25\linewidth}|Y|Y|}
\hline

\rowcolor{softBlue}
\multicolumn{1}{|c|}{\textbf{Logic}}
&
\multicolumn{1}{c|}{\textbf{মূল ধারণা}}
&
\multicolumn{1}{c|}{\textbf{উদাহরণ}}
\\
\hline

Sequence
&
একটির পর একটি কাজ ধারাবাহিকভাবে সম্পন্ন হয়
&
যোগফল, ক্ষেত্রফল, Average
\\
\hline

Selection
&
কোনো Condition অনুযায়ী একটি পথ নির্বাচন করা হয়
&
শর্তভিত্তিক Problem
\\
\hline

Repetition
&
একই কাজ প্রয়োজন অনুযায়ী বারবার করা হয়
&
পুনরাবৃত্তিমূলক Problem
\\
\hline

\end{tabularx}

Selection ও Repetition পরবর্তী Programming অংশে
বিস্তারিত আলোচনা করা হবে।

%================================================
% Characteristics
%================================================
\booksubheading{একটি ভালো Algorithm-এর বৈশিষ্ট্য}

একটি ভালো Algorithm—

\begin{compacttwocolpoints}

\item সহজ ও বোধগম্য হবে।

\item ধাপগুলো সঠিক ও যৌক্তিক ক্রমে থাকবে।

\item প্রতিটি ধাপ স্পষ্ট ও দ্ব্যর্থহীন হবে।

\item অপ্রয়োজনীয় ধাপ থাকবে না।

\item প্রয়োজনীয় Input ও Output পরিষ্কার থাকবে।

\item Problem-এর সঠিক Solution প্রদান করবে।

\item বাস্তবায়নের উপযোগী হবে।

\end{compacttwocolpoints}

%================================================
% Algorithm vs Program
%================================================
\booksubheading{Algorithm ও Program-এর সম্পর্ক}

Algorithm এবং Program একই বিষয় নয়।

\textbf{Algorithm} হলো Problem সমাধানের পরিকল্পিত ধাপ।

\textbf{Program} হলো সেই ধাপগুলোকে Programming Language-এ
Code আকারে বাস্তবায়ন।

\begin{tabularx}{\linewidth}{|>{\bfseries}p{0.22\linewidth}|Y|Y|}
\hline

\rowcolor{softBlue}
\multicolumn{1}{|c|}{\textbf{বিষয়}}
&
\multicolumn{1}{c|}{\textbf{Algorithm}}
&
\multicolumn{1}{c|}{\textbf{Program}}
\\
\hline

রূপ
&
Problem সমাধানের ধাপ
&
Programming Code
\\
\hline

প্রকাশ
&
সাধারণ ভাষা বা সহজ নির্দেশনায় লেখা যায়
&
Programming Language-এ লেখা হয়
\\
\hline

উদ্দেশ্য
&
Solution পরিকল্পনা করা
&
Solution Computer-এ বাস্তবায়ন করা
\\
\hline

অবস্থান
&
সাধারণত Coding-এর আগে
&
Algorithm বা Design-এর পরে
\\
\hline

\end{tabularx}

\begin{center}

\textbf{
Algorithm = Solution-এর পরিকল্পনা
}

\vspace{0.5mm}

\textbf{
Program = সেই পরিকল্পনার Code
}

\end{center}

Algorithm থেকে Program Code-এ যাওয়ার সাধারণ ধারণা—

\begin{center}

\textbf{
Problem
\(\longrightarrow\)
Algorithm
\(\longrightarrow\)
Python Code
\(\longrightarrow\)
Output
}

\end{center}

%================================================
% Common Mistakes
%================================================
\booksubheading{Algorithm লেখার সময় সাধারণ ভুল}

\begin{compacttwocolpoints}

\item ধাপের সঠিক ক্রম না রাখা।

\item Input গ্রহণের প্রয়োজনীয় ধাপ বাদ দেওয়া।

\item Formula বা Calculation ভুল লেখা।

\item Output প্রদর্শনের ধাপ বাদ দেওয়া।

\item অপ্রয়োজনীয় ধাপ যোগ করা।

\item কোনো ধাপ অস্পষ্ট বা দ্ব্যর্থকভাবে লেখা।

\end{compacttwocolpoints}

%================================================
% Final Reminder
%================================================
\begin{rememberbox}

\textbf{\color{deepOrange}মনে রাখবে}

\begin{itemize}

\item \textbf{Algorithm} হলো Problem সমাধানের
\textbf{Step-by-Step পদ্ধতি}।

\item Algorithm সাধারণত Program Coding-এর আগে তৈরি করা হয়।

\item Algorithm-এর সাধারণ কাঠামো—

\begin{center}

\textbf{
শুরু
\(\rightarrow\)
Input
\(\rightarrow\)
Process
\(\rightarrow\)
Output
\(\rightarrow\)
শেষ
}

\end{center}

\item Algorithm নির্দিষ্ট Programming Language-এর
Syntax-এর ওপর নির্ভরশীল নয়।

\item Algorithm-এ
\textbf{Sequence, Selection ও Repetition}
ধরনের Logic থাকতে পারে।

\item Algorithm থেকে Programming Code তৈরি করা সহজ হয়।

\item Algorithm-এর ধাপগুলো
\textbf{স্পষ্ট, সংক্ষিপ্ত ও যৌক্তিক ক্রমে} হওয়া উচিত।

\item Algorithm শব্দটি
\textbf{Musa Al-Khwarizmi}-এর নামের সঙ্গে সম্পর্কিত।

\item দ্রুত মনে রাখবে—

\begin{center}

\textbf{
Problem বুঝি
\(\longrightarrow\)
Solution-এর ধাপ সাজাই
\(\longrightarrow\)
Algorithm লিখি
\(\longrightarrow\)
Code করি
}

\end{center}

\end{itemize}

\end{rememberbox}

\small

%================================================
% 5.24 Flowchart
%================================================

